\documentclass[11pt]{article}
\usepackage[margin=1in]{geometry}
\usepackage{amsmath,amssymb,amsthm,mathtools,microtype,hyperref}
\hypersetup{hidelinks}
\newtheorem{theorem}{Theorem}
\newtheorem{lemma}{Lemma}
\newcommand{\ddc}{dd^c}
\newcommand{\MA}{\operatorname{MA}}
\newcommand{\TV}{\mathrm{TV}}
\newcommand{\E}{\mathbb E}
\title{Shrinking Hermitian potentials and Hopf defect measures}
\author{}
\date{4 October 2026}
\begin{document}
\maketitle
\begin{abstract}
Every defect of an honest fixed-mass shrinking sequence on a compact
Hermitian manifold is a sum of traces of positive closed intermediate
currents annihilating smooth $dd^c$-closed complementary test forms.
Its support lies in the
degeneracy locus of every global closed semipositive $(1,1)$-form.
On the standard Hopf threefold we classify all such defects: they are
exactly the probability-Haar fiber lifts of arbitrary positive measures
on $\mathbb P^2$ of the prescribed excess mass. We construct actual
smooth admissible sequences, with exact mass and any prescribed
shrinking bounds, by controlling mixed quadratic terms, repairing
pullbacks through linear automorphisms and taking diagonal limits.
Possible supports are exactly the nonempty closed fiber-saturated sets.
A lift charges no pluripolar set precisely when its base measure
charges none. Defects may therefore have any fraction of their mass
on a pluripolar elliptic fiber, while every defect is atomless as a
point measure. For every Hermitian metric on the same Hopf manifold,
the complete classification persists, with conditional fiber probability given
by normalized metric area. Exact mass is then obtained by a
quadratic-plus-linear root. Sufficiency on arbitrary Hermitian
manifolds remains open.
\end{abstract}

\section{The setting and the results}
Use $\ddc=(i/(2\pi))\partial\bar\partial$ and
$d^c=(i/(4\pi))(\bar\partial-\partial)$. Let
\[
 X=(\mathbb C^3\setminus\{0\})/\langle z\mapsto e^{-1}z\rangle,
 \quad a=\log|z|^2,
 \quad \alpha=\ddc a,\quad \beta=da\wedge d^c a,
 \quad \omega=\alpha+\beta.
\]
The forms descend, and $\omega$ is a smooth positive Hermitian form.
We use the background identities supplied in S01:
\[
 \ddc\omega=-\alpha^2,\qquad
 \ddc\omega^2=0,\qquad V:=\int_X\omega^3=6,
\]
and the existence of smooth normalized $\omega$-psh functions of
arbitrarily large Monge--Amp\`ere mass. Here $u$ is a genuine smooth
$\omega$-psh function when $\omega+\ddc u\geq0$. Put
\[
 \MA(f)=(\omega+\ddc f)^3,\qquad M(f)=\int_X\MA(f)
\]
for a smooth real function $f$.

An honest shrinking sequence with excess mass $m=b-V>0$ has
\[
 w_j=r_jU_j,\qquad \sup_X U_j=0,\qquad
 \omega+\ddc U_j\geq0,\qquad r_j\longrightarrow0,
 \qquad M(w_j)=b.
\]
Its defect is a weak limit $\MA(w_j)\rightharpoonup\omega^3+\sigma$.
The shrinking results stated in S01 give capacity convergence to zero,
the fixed shallow/deep conclusions and the linear-current conclusions
for precisely these sequences. We preserve all their hypotheses.

\begin{theorem}\label{thm:smooth}
For every $b>V$ and every positive sequence $\varepsilon_j\to0$,
there is an honest smooth shrinking sequence with
$0<r_j<\varepsilon_j$ such that
\[
 M(w_j)=b,\qquad
 \MA(w_j)\rightharpoonup\frac bV\omega^3.
\]
Its defect is the prescribed smooth measure
$\sigma=((b-V)/V)\omega^3$. Its support is $X$, and it charges no
pluripolar set.
\end{theorem}

The group $G=U(3)\times(\mathbb R/\mathbb Z)$ acts by
$(A,t)[z]=[e^tAz]$. For a probability $\rho$ on $G$ and a finite
measure $\sigma$, define the orbital average by continuous tests:
\[
 \mathcal A_\rho\sigma=\int_G g^*\sigma\,d\rho(g),\qquad
 \int_X\varphi\,g^*\sigma=
 \int_X\varphi\circ g^{-1}\,d\sigma.
\]

\begin{theorem}\label{thm:orbit}
Suppose $\sigma$ is already realized by an honest smooth shrinking
sequence on this $(X,\omega)$ with fixed mass $b>V$.
For every probability $\rho$ on $G$, including singular or atomic
probabilities, $\mathcal A_\rho\sigma$ is realized by another such
sequence of exactly the same mass. Its shrinking parameters can satisfy
any prescribed bounds $0<r_j<\varepsilon_j\to0$.
\end{theorem}

This is closure under the orbital convex hull of realizable defects.
Combined with the necessary nonlinear current structure and repaired
linear automorphisms, it yields a complete classification below.

Let $\pi:X\to\mathbb P^2$ be the projectivization, and $h_p$ the
probability Haar measure on the elliptic fiber $E_p=\pi^{-1}(p)$.
For a positive finite base measure $\mu$, its Haar lift $L\mu$ is
defined by
\[
 \int_X\varphi\,d(L\mu)=\int_{\mathbb P^2}
       \left(\int_{E_p}\varphi\,dh_p\right)d\mu(p).
\]

\begin{theorem}\label{thm:classification}
At fixed mass $b>6$, the defects of all honest smooth shrinking
sequences on this Hopf threefold are exactly
\[
 \{L\mu:\mu\geq0,\ \mu(\mathbb P^2)=b-6\}.
\]
Every such measure is realized with any prescribed bounds
$0<r_j<\varepsilon_j\to0$. Possible supports are precisely
$\pi^{-1}(F)$ for nonempty closed $F\subset\mathbb P^2$.
Every defect has no point atoms. The lift $L\mu$ charges no pluripolar
set if and only if $\mu$ charges no pluripolar set on $\mathbb P^2$.
\end{theorem}

\section{The positive quadratic carrier}
\begin{lemma}\label{lem:carrier}
For every smooth real $f$,
\[
 M(f)=V+3Q(f),\qquad
 Q(f)=\int_X\omega\wedge(\ddc f)^2
     =\int_X df\wedge d^c f\wedge\alpha^2\geq0.
\]
Consequently $M(\lambda f)-V=\lambda^2(M(f)-V)$ for real $\lambda$.
If $0<t<1$, $\ddc f\geq-t\omega$ and $m=M(f)-V$, put
$H=\ddc f+t\omega\geq0$ and $\nu=3\omega\wedge H^2$. Then
\[
 \int_X\omega^2H=tV,\qquad
 \int_X\omega H^2=m/3+t^2V,\qquad
 \int_XH^3=tm+t^3V
\]
and
\[
 \|\MA(f)-\omega^3-\nu\|_{\TV}
       \leq 6tV+2tm+4t^3V.
\]
Thus the complete defect limit of an honest fixed-mass sequence
equals the weak limit of its positive quadratic carriers.
\end{lemma}

\begin{proof}
Stokes and $\ddc\omega^2=0$ give
$\int\omega^2\ddc f=0$. Also
$\int(\ddc f)^3=0$, because this smooth form is exact. Expanding
$M(f)$ proves the first equality. As
$\ddc(f\ddc f)=(\ddc f)^2$, integration by parts gives
$Q(f)=\int f\ddc f\wedge\ddc\omega$. The identity
\[
 \ddc(f^2/2)=df\wedge d^c f+f\ddc f
\]
and closedness of $\ddc\omega$ yield
$Q(f)=-\int df\wedge d^c f\wedge\ddc\omega$.
Substitute $\ddc\omega=-\alpha^2$. Positivity follows from
$df\wedge d^c f\geq0$ and $\alpha\geq0$.
Expanding the three integrals involving $H$ gives the stated masses.

The actual positive-form expansion is
\[
 \MA(f)=(1-t)^3\omega^3+3(1-t)^2\omega^2H
                     +(1-t)\nu+H^3.
\]
All four terms are positive. The triangle inequality therefore bounds
the total variation after subtracting $\omega^3+\nu$ by
\[
 [1-(1-t)^3]V+3(1-t)^2tV+t(m+3t^2V)+(tm+t^3V),
\]
which is at most the claimed expression. This is a measure estimate,
rather than merely an equality of integrated masses.
\end{proof}

\section{Averaging actual potentials}
The group action is holomorphic, transitive and preserves $\omega$:
unitaries preserve $a$, while real scaling adds a constant to it.
There is an invariant probability $h$ on $G$. One explicit construction
takes a complex $3\times3$ Gaussian matrix $Z$, sets
$A=Z(Z^*Z)^{-1/2}$ almost surely and independently chooses $t$
uniformly modulo one. The Gaussian matrix is invertible almost surely.
Its left and right unitary invariance and the equivariance of this
polar expression give left and right invariance of $h$.

For a smooth $q$, transitivity and invariance make
$\int_G q\circ g\,dh(g)$ constant. Differentiating the smooth average
and using preservation of $\omega$ gives
\[
 \int_Gg^*(\ddc q+s\omega)\,dh(g)=s\omega.
 \tag{1}\label{eq:form-average}
\]
For every finite positive measure $\eta$,
\[
 \int_Gg^*\eta\,dh(g)=\frac{\eta(X)}V\omega^3.
 \tag{2}\label{eq:measure-average}
\]
Indeed the average of $\varphi\circ g^{-1}$ is constant in the point;
integrating that constant against the invariant probability
$\omega^3/V$ identifies it as $V^{-1}\int\varphi\omega^3$.

Let $f=su$ be an honest fixed-mass shrink with $M(f)=V+m$, and put
$H=\ddc f+s\omega$, $\nu=3\omega H^2$.
For any $N$ translates define the actual smooth potential and form
\[
 W=N^{-1/2}\sum_{l=1}^N f\circ g_l,\qquad
 \tau=s\sqrt N,\qquad
 K=\ddc W+\tau\omega=N^{-1/2}\sum_{l=1}^Ng_l^*H\geq0.
\]
Their quadratic carrier decomposes into positive measures
\[
 3\omega K^2=D+C,\qquad
 D=\frac1N\sum_lg_l^*\nu,\qquad
 C=\frac6N\sum_{l<a}\omega\wedge g_l^*H\wedge g_a^*H.
 \tag{3}\label{eq:diagonal}
\]
In particular nonlinear products are not assumed to commute with
averaging. We will estimate the actual mixed term $C$.

Suppose the $g_l$ are independent with law $\rho$ having density
at most $L<\infty$ relative to $h$. Positivity and
\eqref{eq:form-average} show
\[
 0\leq B:=\int_Gg^*H\,d\rho(g)\leq Ls\omega.
\]
Independence and Fubini imply
\[
 \E[C(X)]=3(N-1)\int_X\omega B^2
                    \leq3(N-1)L^2s^2V.
 \tag{4}\label{eq:cross}
\]
For $\rho=h$ the expectation is exactly $3(N-1)s^2V$.
For a real test $\varphi$ with $\|\varphi\|_\infty\leq1$, the
independent variables $\int\varphi\,g_l^*\nu$ have absolute value
at most $m+3s^2V$. Their empirical mean $\int\varphi D$ has expectation
$\int\varphi\,\mathcal A_\rho\nu$ and variance at most
\[
 \frac{(m+3s^2V)^2}{N}.
 \tag{5}\label{eq:variance}
\]

\section{The smooth prescribed defect}
\begin{proof}[Proof of Theorem~\ref{thm:smooth}]
Fix $m=b-V>0$ and a countable uniformly dense collection
$(\varphi_k)$ in the real unit ball of $C(X)$. Choose
$\delta_j\downarrow0$ with $\delta_j<m/10$. First choose
$N_j\to\infty$ so large that
\[
 \frac{j(m+3V)^2}{N_j\delta_j^2}<\frac14.
\]
Then choose a smooth normalized $\omega$-psh $u_j$ of sufficiently
large mass and set
\[
 s_j=\sqrt{\frac{m}{M(u_j)-V}},\qquad f_j=s_ju_j.
\]
Lemma~\ref{lem:carrier} gives $M(f_j)=b$. Large source mass makes
$s_j$ arbitrarily small, so we impose
\[
 s_j\sqrt{N_j}<\min(\varepsilon_j/4,1/j,1/4),\qquad
 3VN_js_j^2<m/2,\qquad
 3(N_j-1)s_j^2V<\delta_j/4.
\]

Sample $N_j$ independent translates with law $h$. Chebyshev's
inequality, \eqref{eq:variance} and a union bound show that the
probability of an empirical error exceeding $\delta_j$ on any of
the first $j$ tests is less than $1/4$. Markov's inequality and
\eqref{eq:cross} show that the probability of $C(X)>\delta_j$ is
less than $1/4$. Thus a deterministic tuple satisfies both bounds.
Fix one such tuple at each step. Equations~\eqref{eq:measure-average}
and \eqref{eq:diagonal} imply
\[
 D_j\rightharpoonup\frac mV\omega^3,\qquad
 \|C_j\|_{\TV}=C_j(X)\leq\delta_j\longrightarrow0,
 \qquad 3\omega K_j^2\rightharpoonup\frac mV\omega^3.
\]
The dense tests suffice because all these measure masses are bounded.

We now fix the mass exactly. With $\tau_j=s_j\sqrt{N_j}$,
Lemma~\ref{lem:carrier} and $D_j(X)=m+3s_j^2V$ give
\[
 d_j:=M(W_j)-V=m+C_j(X)-3(N_j-1)s_j^2V\longrightarrow m,
 \qquad d_j>m/2.
\]
Set $\lambda_j=\sqrt{m/d_j}\to1$, so $\lambda_j\leq\sqrt2$.
Let
\[
 a_j=\frac1{N_j}\sum_lu_j\circ g_l,\qquad
 U_j=a_j-\sup_Xa_j,\qquad r_j=\lambda_j\tau_j,
 \qquad w_j=r_jU_j.
\]
The normalized $U_j$ is smooth and $\omega$-psh, and
$0<r_j<\varepsilon_j$. Since $w_j$ differs from
$\lambda_jW_j$ by a constant, the exact quadratic scaling gives
\[
 M(w_j)=V+\lambda_j^2d_j=b,\qquad
 \omega+\ddc w_j\geq(1-r_j)\omega,
 \qquad \ddc w_j+r_j\omega=\lambda_jK_j.
\]
Its quadratic carrier tends to $(m/V)\omega^3$. The total-variation
estimate in Lemma~\ref{lem:carrier} proves
$\MA(w_j)\rightharpoonup(b/V)\omega^3$ against every continuous test.
All potentials and products are actual smooth objects at each step.

This is exactly an honest shrinking sequence, so the supplied capacity,
shallow/deep and linear-current conclusions apply. Its defect has full
support because $\omega$ is positive. Locally a pluripolar set lies in
the $-\infty$ locus of a non-identically-$-\infty$ psh function.
Local integrability \cite[I.4.17 and I.5.3]{Demailly} makes that
locus Lebesgue-null, and hence null
for the smooth measure $(m/V)\omega^3$.
\end{proof}

\section{Arbitrary orbital laws}
\begin{proof}[Proof of Theorem~\ref{thm:orbit}]
Write the source sequence as $f_k=s_ku_k$, and put
$H_k=\ddc f_k+s_k\omega$, $\nu_k=3\omega H_k^2$.
Lemma~\ref{lem:carrier} yields
\[
 \nu_k\rightharpoonup\sigma,\qquad
 \nu_k(X)=m+3s_k^2V,
 \qquad\int_Gg^*H_k\,dh(g)=s_k\omega.
\]
Approximate $\rho$ weakly by probabilities $\rho_j$ with finite
density bounds $L_j$ relative to $h$. Explicitly, restrict $h$
to a shrinking neighborhood $B_j$ of the identity, normalize by
$h(B_j)$, and convolve with $\rho$. Nonempty neighborhoods have
positive invariant mass: finitely many translates cover the compact
group, so zero mass would contradict $h(G)=1$. The resulting density
is at most $L_j=1/h(B_j)$. Uniform continuity on the compact group
shows that these convolutions tend weakly to $\rho$.

Choose $\delta_j\downarrow0$ and $N_j\to\infty$ to meet the same
variance/union bound as above. Choose increasing source indices $k_j$
so large that $s_{k_j}<1$ and
\[
 s_{k_j}\sqrt{N_j}<\min(\varepsilon_j/4,1/j,1/4),\qquad
 3VN_js_{k_j}^2<m/2,\qquad
 3(N_j-1)L_j^2s_{k_j}^2V<\delta_j/4.
\]
The limits $s_k\to0$ make all these conditions possible.
Sample with law $\rho_j$ and select deterministic tuples using
\eqref{eq:cross}--\eqref{eq:variance}. Their diagonal measures differ
from $\mathcal A_{\rho_j}\nu_{k_j}$ by at most $\delta_j$ on the
first $j$ dense tests, and their mixed measures have mass at most
$\delta_j$.

For completeness, the expectation measures have the desired limit.
For fixed continuous $\varphi$, the family
$\{\varphi\circ g^{-1}:g\in G\}$ is compact in the uniform topology.
Take a finite uniform net of this family. Weak convergence at its
finitely many functions, together with uniformly bounded masses,
proves
\[
 \sup_{g\in G}\left|
      \int_X\varphi\,g^*\nu_k-\int_X\varphi\,g^*\sigma
                  \right|\longrightarrow0.
\]
Also $g\mapsto\int\varphi\,g^*\sigma$ is continuous. Hence
$\mathcal A_{\rho_j}\nu_{k_j}\rightharpoonup\mathcal A_\rho\sigma$.
The selected tuples give
\[
 D_j\rightharpoonup\mathcal A_\rho\sigma,\qquad
 C_j\longrightarrow0\ \text{in total variation},\qquad
 3\omega K_j^2\rightharpoonup\mathcal A_\rho\sigma.
\]

Repeat the exact normalization from the preceding proof, now with
$s_j=s_{k_j}$. The same formula gives
$d_j=m+C_j(X)-3(N_j-1)s_j^2V>m/2$ and $d_j\to m$.
Set $\lambda_j=\sqrt{m/d_j}$, normalize the average of the
$u_{k_j}\circ g_l$ to supremum zero and take
$r_j=\lambda_js_j\sqrt{N_j}<\varepsilon_j$.
The resulting actual smooth $w_j=r_jU_j$ has $M(w_j)=b$ and
quadratic carrier $\lambda_j^2(D_j+C_j)$ tending to
$\mathcal A_\rho\sigma$. Lemma~\ref{lem:carrier} gives the complete
weak limit $\MA(w_j)\rightharpoonup\omega^3+\mathcal A_\rho\sigma$.
The honest shrinking hypotheses, including their capacity and
fixed-level consequences, are retained.
\end{proof}

Finite laws realize every convex combination
$\sum_a\theta_ag_a^*\sigma$ with $\theta_a\geq0$ and
$\sum_a\theta_a=1$. Its support is the union of the supports of
the positively weighted translates. Taking
$\rho=\theta\delta_e+(1-\theta)h$ realizes
\[
 \theta\sigma+(1-\theta)\frac mV\omega^3,
 \qquad 0\leq\theta\leq1.
\]
For $\theta<1$ its support is $X$. If the source $\sigma$ charges
a pluripolar set $P$ and $\theta>0$, the mixture charges $P$ with
mass $\theta\sigma(P)$. The fiber constructions below supply such
charging sources.

\section{Necessary currents on arbitrary Hermitian manifolds}
In this section $(Y,h)$ is any connected compact Hermitian manifold
of dimension $n$, and $v_j=s_ju_j$ is an honest shrinking sequence
with fixed mass $b$ and defect $\sigma$:
\[
 \sup_Yu_j=0,\quad h+\ddc u_j\geq0,\quad s_j\to0,\quad
 \int_Y(h+\ddc v_j)^n=b,\quad
 (h+\ddc v_j)^n\rightharpoonup h^n+\sigma.
\]

\begin{theorem}\label{thm:general}
After a common subsequence there are positive $d$-closed currents
$K_k$ of bidegree $(k,k)$, $2\leq k\leq n-1$, such that
\[
 \sigma=\sum_{k=2}^{n-1}\binom nk h^{n-k}\wedge K_k.
 \tag{6}\label{eq:general-defect}
\]
Each $K_k$ has zero pairing with every smooth $d$-closed real form
of complementary degree and every smooth $dd^c$-closed form of
bidegree $(n-k,n-k)$. The first and top positive powers vanish,
the top one in total variation. For every global smooth closed
semipositive real $(1,1)$-form $\vartheta$,
\[
\operatorname{supp}\sigma\subset\{\vartheta^n=0\}.
\]
In dimension three the same support conclusion holds for every
global smooth semipositive pluriclosed $(1,1)$-form $\vartheta$.
These conditions are necessary; their sufficiency is not asserted.
\end{theorem}

\begin{proof}
Put $H_j=\ddc v_j+s_jh=s_j(h+\ddc u_j)\geq0$.
Discard finitely many indices so that $s_j<1/2$. Positivity in the
actual expansion
\[
 (h+\ddc v_j)^n=\sum_{k=0}^n\binom nk
                  (1-s_j)^{n-k}h^{n-k}H_j^k
\]
and the fixed mass give
\[
 \int_Yh^{n-k}H_j^k\leq\frac{2^{n-k}b}{\binom nk}.
 \tag{7}\label{eq:power-bounds}
\]
In finitely many charts these positive traces bound the coefficients;
positivity also bounds off-diagonal coefficients by diagonal ones.
Weak compactness of bounded coefficient measures gives simultaneous
positive current limits $H_j^k\rightharpoonup K_k$.

The first power has mass $O(s_j)$. Use a smooth Gauduchon metric
$\gamma$ from the supplied trace statement, so
$\ddc\gamma^{n-1}=0$. Smooth Stokes gives
\[
 \int_Y\gamma^{n-1}H_j=s_j\int_Y\gamma^{n-1}h.
\]
Metric comparison bounds the $h$-trace by this one. Thus $K_1=0$.
Since $dH_j=s_jdh$,
\[
 d(H_j^k)=ks_jdh\wedge H_j^{k-1}\longrightarrow0
\]
against every fixed smooth test, by \eqref{eq:power-bounds}.
Hence $dK_k=0$.

Let $\eta$ be a smooth $d$-closed real complementary form.
The smooth exactness
$(\ddc v_j)^k=d(d^cv_j\wedge(\ddc v_j)^{k-1})$ yields
$\int\eta(\ddc v_j)^k=0$. Expanding $\ddc v_j=H_j-s_jh$ gives
\[
 \int_Y\eta H_j^k
 =-\sum_{r=0}^{k-1}\binom kr(-s_j)^{k-r}
                 \int_Y\eta\wedge h^{k-r}H_j^r=O(s_j).
\]
Indeed the absolute value of each integral is bounded by a fixed
multiple of the corresponding trace in \eqref{eq:power-bounds}.
Passing to limits proves the asserted zero pairings. Taking $k=n$
and $\eta=1$ shows $\int H_j^n=O(s_j)$; positivity gives total-variation
convergence to zero and $K_n=0$. Passing to limits in the actual
positive expansion now proves \eqref{eq:general-defect}.

For the stronger test-form assertion, take a smooth $(n-k,n-k)$-form
$\eta$ with $\ddc\eta=0$. The exact smooth identity
\[
 (\ddc v_j)^k=\ddc\bigl(v_j(\ddc v_j)^{k-1}\bigr)
\]
and integration by parts again give $\int\eta(\ddc v_j)^k=0$.
The same binomial expansion and trace bounds give
$\int\eta H_j^k=O(s_j)$, and hence $\int\eta K_k=0$.

For a global closed semipositive $\vartheta$, take
$\eta=\vartheta^{n-k}$. The positive measure
$\vartheta^{n-k}\wedge K_k$ has total mass zero and thus is zero.
Where $\vartheta$ is strictly positive, its power dominates a positive
multiple of $h^{n-k}$ on every relatively compact subregion. Therefore
each $K_k$ vanishes there, proving the support assertion.
The form must be global: a local K\"ahler form cannot be used in
the closed-form pairing argument.

In dimension three use any global $\vartheta\geq0$ with
$\ddc\vartheta=0$ as the stronger complementary test. Its positive
pairing with $K_2$ has zero mass, so $K_2$ vanishes where
$\vartheta$ is strictly positive. This proves the pluriclosed
support assertion. More generally a global $dd^c$-closed complementary
form whose wedge with positive currents is nonnegative eliminates
its corresponding $K_k$ wherever it dominates a positive multiple
of $h^{n-k}$.
\end{proof}

For $n=3$, the formula is $\sigma=3h\wedge K_2$.
For $n\leq2$ the sum is empty, so these honest shrinking hypotheses
cannot hold with positive excess mass. No cohomological
representation theorem is required: the pairings themselves are
the precise assertion. In particular a positive pluriclosed metric
on a threefold excludes these positive fixed-mass defects. We do not
infer $dd^c$-exactness of a limiting current merely from limits of
$dd^c$-exact smooth forms without a closed-range argument.

\section{Every Hopf defect has Haar fibers}
The Hopf fibration is locally a holomorphic product $U\times E$,
with $E=\mathbb C/(\mathbb Z+2\pi i\mathbb Z)$ and fiber coordinate
$w=\log z_0$ in a base chart. Transition functions translate $w$
by a holomorphic base function. Thus fiber Haar probability is
canonical. Its average of a continuous function is continuous in the
base, by a product chart and uniform continuity. Consequently $L$
preserves weak convergence of finite base measures.

\begin{lemma}\label{lem:necessity}
Every Hopf defect has the form $\sigma=L\mu$ with
$\mu=\pi_*\sigma\geq0$ of mass $m=b-6$. Along the sequence on which
the defect limit is fixed,
\[
 (\ddc w_j+s_j\omega)^2\rightharpoonup\frac16\pi^*\mu.
\]
Moreover $\operatorname{supp}L\mu=\pi^{-1}(\operatorname{supp}\mu)$,
and $L\mu$ has no point atoms.
\end{lemma}

\begin{proof}
Theorem~\ref{thm:general} gives a positive closed $(2,2)$-current
$K$ with $\sigma=3\omega K$. Its pairing with
$\alpha=\pi^*\omega_{\mathrm{FS}}$ is zero. Since both are positive,
$\alpha\wedge K=0$ as a measure.

In a product chart with base coordinates $x_1,x_2$, the positive
Hermitian coefficient matrix of $K$ is indexed by
\[
 dx_1\wedge dx_2,\qquad dx_1\wedge dw,\qquad dx_2\wedge dw.
\]
This can be interpreted rigorously using Radon--Nikodym densities
relative to the sum of its diagonal coefficient measures.
The form $\alpha$ is strictly positive horizontally. Thus
$\alpha K=0$ forces the two diagonal entries containing $dw$ to
vanish. The positive-matrix inequality
$|a_{ij}|^2\leq a_{ii}a_{jj}$ forces their cross entries to vanish too.
Only the horizontal top-base coefficient remains.

Closedness says that both real vertical distributional derivatives
of that coefficient are zero. Test it against a compactly supported
base function times a Fourier mode
$e^{2\pi ik\operatorname{Re}w+i\ell\operatorname{Im}w}$.
For every nonconstant mode at least one derivative kills its
coefficient. Density of trigonometric polynomials then shows that
the scalar measure is a product of a base measure with constant
fiber area. The translation transition functions preserve this
description, so globally
\[
 K=\pi^*\tau
\]
for a positive top-degree base current, equivalently a positive
measure $\tau$. This is pullback through a submersion, defined by
the local horizontal product construction.

On the fiber,
\[
 \beta|_E=\frac1\pi\,d(\operatorname{Re}w)
                      \wedge d(\operatorname{Im}w),
 \qquad \int_E\beta=2.
\]
Hence $L\mu=(1/2)\beta\wedge\pi^*\mu$. Horizontal components of
$\beta$ disappear in this wedge, and $\alpha\pi^*\tau=0$.
Therefore
\[
 \sigma=3\beta\wedge\pi^*\tau=L(6\tau).
\]
Set $\mu=6\tau=\pi_*\sigma$; its mass is $m$.
Once $\sigma$ is fixed, every subsequential current limit is the
same $K=\pi^*\mu/6$. The trace bounds give precompactness, so the
entire carrier-current sequence converges to this $K$.
Locally the lift is a product with full-support atomless Haar
probability, proving the support and atom assertions.
\end{proof}

\section{Prescribing a pluripolar fiber}
\begin{lemma}\label{lem:fiber}
For every $p\in\mathbb P^2$, the measure $mh_p$ is an honest
shrinking defect with exactly mass $b=6+m$ and any prescribed
shrinking bounds. The fiber $E_p$ is globally $\omega$-pluripolar.
\end{lemma}

\begin{proof}
Take the diffuse sequence of Theorem~\ref{thm:smooth}, writing
$f_k=s_ku_k$ and $H_k=\ddc f_k+s_k\omega$.
Since $\omega^3=3\alpha^2\beta$ and $\int_E\beta=2$, the base
normalization is $\int_{\mathbb P^2}\omega_{\mathrm{FS}}^2=1$.
Its defect $(m/6)\omega^3=L(m\omega_{\mathrm{FS}}^2)$ and
Lemma~\ref{lem:necessity} give the indispensable current limit
\[
 H_k^2\rightharpoonup(m/6)\alpha^2,
 \qquad \operatorname{mass}(H_k)=O(s_k).
\]

Every $A\in GL(3,\mathbb C)$ commutes with the deck dilation and
defines an automorphism of $X$ and a projective automorphism $a$
of the base. Fix $A$. Choose $C_A\geq1$ with
$A^*\omega\leq C_A\omega$, and set
$R_A=C_A\omega-A^*\omega\geq0$.
For $f_k^A=f_k\circ A$ and $t_k=C_As_k$ the repaired positive
carrier is
\[
 H_k^A=\ddc f_k^A+t_k\omega=A^*H_k+s_kR_A\geq0.
\]
The difference
\[
 (H_k^A)^2-A^*(H_k^2)
        =2s_kA^*H_k\wedge R_A+s_k^2R_A^2
\]
has current mass $O_A(s_k^2)$: fixed-metric comparison gives
$\operatorname{mass}(A^*H_k)=O_A(s_k)$.
It follows that
\[
 (H_k^A)^2\rightharpoonup(m/6)\pi^*(a^*\omega_{\mathrm{FS}}^2),
 \quad
 3\omega(H_k^A)^2\rightharpoonup L(m a^*\omega_{\mathrm{FS}}^2).
\]
The latter measure has mass $m$, by change of variables on the base.

Let $d_k^A=M(f_k^A)-6$. Lemma~\ref{lem:carrier} gives
$d_k^A=3\int\omega(H_k^A)^2-3t_k^2V\to m$.
For sufficiently late $k$, take
$\lambda_k^A=\sqrt{m/d_k^A}\to1$. Exact scaling gives
$M(\lambda_k^Af_k^A)=b$. The underlying function
\[
 u_k^A=(u_k\circ A)/C_A
\]
is smooth, has supremum zero, and satisfies
$\omega+\ddc u_k^A\geq\omega-C_A^{-1}A^*\omega\geq0$.
Thus $\lambda_k^Af_k^A=r_k^Au_k^A$ with
$r_k^A=\lambda_k^AC_As_k\to0$. Its carrier is
$\lambda_k^AH_k^A$. The total-variation reduction proves the complete
limit $\MA(\lambda_k^Af_k^A)\rightharpoonup
\omega^3+L(m a^*\omega_{\mathrm{FS}}^2)$.
The background stays $\omega$ throughout; the repair proves actual
admissibility rather than assuming a pullback formula for these
nonlinear measures.

For $p_0=[1:0:0]$, take
$A_j=\operatorname{diag}(R_j^{-1},1,1)$ with $R_j\to\infty$.
Since
$a_j^*\omega_{\mathrm{FS}}^2=(a_j^{-1})_*\omega_{\mathrm{FS}}^2$
and
\[
 a_j^{-1}[q_0:q_1:q_2]=[R_jq_0:q_1:q_2]\longrightarrow p_0
       \quad(q_0\ne0),
\]
dominated convergence on continuous tests gives
$a_j^*\omega_{\mathrm{FS}}^2\rightharpoonup\delta_{p_0}$.
The exceptional hyperplane is null for smooth base volume.
Weak continuity of $L$ yields
$L(m a_j^*\omega_{\mathrm{FS}}^2)\rightharpoonup mh_{p_0}$.

For each fixed $A_j$, choose a sufficiently late actual member of
its just-constructed honest sequence. Require parameter less than
$\varepsilon_j$ and error at most $1/j$ on the first $j$ functions
of a fixed dense continuous test collection. The mass is already
exactly $b$. These diagonal choices give the complete weak limit
$\omega^3+mh_{p_0}$, with all the original shrinking constraints.
A fixed unitary rotation gives any other $p$.

Finally
\[
 \psi([z])=\log\frac{|z_1|^2+|z_2|^2}{|z|^2}
\]
is well-defined on $X$, locally a genuine psh logarithm minus a
smooth function, and $\ddc\psi\geq-\alpha$.
Thus $\omega+\ddc\psi\geq\beta\geq0$, and its pole set is exactly
$E_{p_0}$. This proves global $\omega$-pluripolarity, also for
unitary images $E_p$. The realized defect assigns its full mass
$m$ to that fiber.
\end{proof}

\section{Classification, supports and pluripolar charging}
\begin{proof}[Proof of Theorem~\ref{thm:classification}]
Necessity is Lemma~\ref{lem:necessity}. For finite atomic
\[
 \mu=m\sum_{a=1}^N\theta_a\delta_{p_a},\qquad
 \theta_a\geq0,\quad\sum_a\theta_a=1,
\]
start with the fiber realization $mh_{p_0}$ of Lemma~\ref{lem:fiber}.
Choose unitary $g_a$ with $g_a^{-1}(p_0)=p_a$ on the base.
Their maps between fibers are translations and preserve Haar
probability. Theorem~\ref{thm:orbit} with the finite law
$\sum_a\theta_a\delta_{g_a}$ realizes
$\sum_a\theta_amh_{p_a}=L\mu$, by actual smooth potentials of
exact mass and arbitrarily small parameters.

For arbitrary $\mu$, take finite Borel partitions of $\mathbb P^2$
whose diameters tend to zero. Replace the mass of each element by
an atom in that element, obtaining $\mu_j\rightharpoonup\mu$ with
$\mu_j(\mathbb P^2)=m$. From each honest realization of $L\mu_j$
choose a late member with parameter $r_j<\varepsilon_j$ and error
at most $1/j$ on the first $j$ dense tests. The selected potentials
are smooth and normalized, their mass remains exactly $b$, and weak
continuity of $L$ gives the full limit
\[
 \MA(r_jU_j)\rightharpoonup\omega^3+L\mu.
\]
Their positivity is $\omega+\ddc(r_jU_j)\geq(1-r_j)\omega$.
The supplied capacity, shallow/deep and linear-current conclusions
apply to these honest shrinking sequences.

By Lemma~\ref{lem:necessity}, supports are exactly preimages of base
supports. Every nonempty closed $F\subset\mathbb P^2$ supports a
measure of mass $m$ with support $F$: take a countable dense sequence
$(p_l)$ in $F$ and the measure $\sum_{l\geq1}m2^{-l}\delta_{p_l}$.
Its realization proves the claimed complete support description.
Atomlessness also follows from the lemma.

It remains to prove the pluripolarity equivalence. If $B$ is a base
pluripolar set, pull back its local nontrivial psh pole witnesses
through the holomorphic submersion. They remain nontrivial, and
$\pi^{-1}(B)$ is pluripolar. Since
$L\mu(\pi^{-1}(B))=\mu(B)$, a non-pluripolar-charging lift forces
the same property for $\mu$.

Conversely suppose $\mu$ charges no base pluripolar set, and let
$P\subset X$ be a measurable pluripolar set. Use countably many
relatively compact product charts $U\times D$ with local nontrivial
psh functions $q(x,w)$ whose pole loci contain the corresponding
pieces of $P$. Smaller product charts still cover $X$. On a fixed
smaller fiber disc $D'$, each restriction $q(x,\cdot)$ is either
nontrivial subharmonic, with area-null pole set, or identically
$-\infty$. The exceptional parameters in the second case are
base-pluripolar. Indeed
\[
 F(x)=\frac1{\operatorname{area}(D')}
                   \int_{D'}q(x,w)\,dA(w)
\]
is a genuine nontrivial base psh function. Joint local integrability
and Fubini make it finite almost everywhere. A common local upper
bound and reverse Fatou prove upper semicontinuity; integrating the
base submean inequalities proves the psh inequalities. If the fiber
restriction is identically $-\infty$, then $F(x)=-\infty$.

Outside the countable union of these base pole sets, every piece of
$P$ has fiber Haar measure zero. The union is $\mu$-null by
hypothesis, and the defining fiber integral shows $L\mu(P)=0$.
Using the local Borel pole supersets proves the same forward assertion
with outer measures. For the reverse assertion, let a Borel null set
$C$ contain $\pi^{-1}(B)$, where $B$ is a possibly nonmeasurable base
pluripolar set. The fiber kernel $p\mapsto h_p(C)$ is Borel: this holds
first for continuous tests by the product trivializations, then for
Borel tests by the monotone-class argument. It equals one on $B$ and
has $\mu$-integral $L\mu(C)=0$, so $\mu^*(B)=0$.
An outer-null set has such a Borel null envelope. This proves the
outer-measure equivalence as well.
\end{proof}

For each $p$ and $0\leq\theta\leq1$, the classification realizes
\[
 \mu=m\theta\delta_p+m(1-\theta)\omega_{\mathrm{FS}}^2,
 \qquad
 \sigma=m\theta h_p+\frac{m(1-\theta)}6\omega^3.
\]
The globally pluripolar fiber $E_p$ carries exactly $m\theta$.
The smooth part charges no pluripolar set. Thus every fraction of
the excess mass can be a pluripolar part, including full-support
charging examples for $0<\theta<1$, a smooth noncharging example
for $\theta=0$ and a fully fiber-supported example for $\theta=1$.
All remain atomless as point measures.

\section{Every Hermitian metric on the same Hopf manifold}
Let $\Omega$ now be any smooth Hermitian metric on this scalar
Hopf threefold. Put $V_\Omega=\int_X\Omega^3$ and
\[
 a_\Omega(p)=\int_{E_p}\Omega>0,\qquad
 h_p^\Omega=\Omega|_{E_p}/a_\Omega(p).
\]
Fiber integration in local smooth bundle trivializations shows that
$a_\Omega$ is smooth and bounded above and below by positive constants.
Define the weighted lift $L_\Omega\mu$ by averaging against
$h_p^\Omega$ in each fiber. These fiber averages are continuous,
so the weighted lift is weakly continuous.

\begin{theorem}\label{thm:arbitrary-metric}
For every $b>V_\Omega$, the defects of all honest smooth
$\Omega$-psh shrinking sequences of exact mass $b$ are precisely
\[
 \{L_\Omega\mu:\mu\geq0,\ \mu(\mathbb P^2)=b-V_\Omega\}.
\]
Every such measure is realized with any prescribed shrinking bounds.
The possible supports are exactly nonempty closed fiber-saturated
sets, defects have no point atoms, and $L_\Omega\mu$ charges no
pluripolar set if and only if $\mu$ charges none.
\end{theorem}

\begin{proof}
Set $m=b-V_\Omega>0$. For necessity apply
Theorem~\ref{thm:general} with $h=\Omega$.
The carrier is positive and closed and annihilates the closed
semipositive $\alpha$. The horizontal-coefficient and vertical-Fourier
argument of Lemma~\ref{lem:necessity} gives $K=\pi^*\tau$, independent
of the background metric. Hence
\[
 \sigma=3\Omega\wedge\pi^*\tau=L_\Omega\mu,
 \qquad \mu=3a_\Omega\tau=\pi_*\sigma,
 \qquad\mu(\mathbb P^2)=m.
\]

Conversely let $\mu$ have mass $m$, and set
$\tau=\mu/(3a_\Omega)$ by multiplication with that positive
continuous function. Put $m_0=6\tau(\mathbb P^2)$ and $\mu_0=6\tau$.
The standard-metric classification realizes $L\mu_0$ with source
mass $6+m_0$. Write its smooth functions as $f_j=s_ju_j$.
Lemma~\ref{lem:necessity} gives
\[
 H_j=\ddc f_j+s_j\omega\geq0,\qquad
 H_j^2\rightharpoonup\pi^*\tau,\qquad
 \operatorname{mass}(H_j)=O(s_j).
\]
Choose $C\geq1$ with $\omega\leq C\Omega$, and put
$R=C\Omega-\omega\geq0$, $t_j=Cs_j$.
The repaired carrier
\[
 G_j=\ddc f_j+t_j\Omega=H_j+s_jR\geq0
\]
has mass $O(s_j)$. Its square differs from $H_j^2$ by
$2s_jH_jR+s_j^2R^2$, a positive current of mass $O(s_j^2)$.
Therefore $G_j^2\rightharpoonup\pi^*\tau$ and
$3\Omega G_j^2\rightharpoonup L_\Omega\mu$.

The exact mass adjustment includes a linear term. Define
\[
 \ell_j=3\int_X\Omega^2\ddc f_j,
 \qquad q_j=3\int_X\Omega(\ddc f_j)^2.
\]
Smooth Stokes eliminates the cubic exact integral, giving
\[
 M_\Omega(\lambda f_j):=\int_X(\Omega+\ddc(\lambda f_j))^3
       =V_\Omega+\lambda\ell_j+\lambda^2q_j.
\]
The carrier estimates give
\[
 \ell_j=3\int\Omega^2G_j-3t_jV_\Omega=O(s_j),
\]
\[
 q_j=3\int\Omega G_j^2-6t_j\int\Omega^2G_j
                               +3t_j^2V_\Omega\longrightarrow m.
\]
For sufficiently late $j$, $q_j>0$. The positive root
\[
 \lambda_j=\frac{-\ell_j+\sqrt{\ell_j^2+4mq_j}}{2q_j}
             \longrightarrow1
\]
satisfies $\lambda_j\ell_j+\lambda_j^2q_j=m$.
Thus $M_\Omega(\lambda_jf_j)=b$ exactly.
The function $U_j=u_j/C$ has supremum zero and
$\Omega+\ddc U_j\geq\Omega-\omega/C\geq0$.
Set $r_j=\lambda_jCs_j$ and $w_j=r_jU_j=\lambda_jf_j$.
Subsequence selection gives $0<r_j<\varepsilon_j$ for every
prescribed sequence of positive bounds. All these functions are
smooth genuine normalized $\Omega$-psh underlying potentials.

The positive carrier of $w_j$ is $\lambda_jG_j$.
Its square tends to $\pi^*\tau$ and its first-power trace tends
to zero. Also
\[
 \int_XG_j^3=t_jq_j+t_j^2\ell_j+t_j^3V_\Omega=O(s_j).
\]
This follows by smooth expansion and Stokes, and positivity makes
it total-variation convergence to zero. The actual positive expansion
is
\[
 \begin{split}
 (\Omega+\ddc w_j)^3={}&(1-r_j)^3\Omega^3
       +3(1-r_j)^2\lambda_j\Omega^2G_j\\
       &+3(1-r_j)\lambda_j^2\Omega G_j^2+\lambda_j^3G_j^3.
 \end{split}
\]
The first term tends to $\Omega^3$, the second and fourth vanish
in total variation, and the third tends weakly to $L_\Omega\mu$.
This proves the complete local weak limit, with exact mass at every
step. The honest shrinking class retains the supplied capacity and
fixed-level conclusions.

The metric has infinite smooth upper mass: choose $c>0$ with
$\Omega\geq c\omega$. For any smooth $\omega$-psh $u$,
\[
 \Omega+\ddc(cu)=c(\omega+\ddc u)+(\Omega-c\omega)
                           \geq c(\omega+\ddc u),
\]
so $M_\Omega(cu)\geq c^3M_\omega(u)$.
The same $c$ times the logarithmic witness in
Lemma~\ref{lem:fiber} makes each $E_p$ globally
$\Omega$-pluripolar, since $\Omega+\ddc(c\psi)\geq
\Omega-c\alpha\geq0$.
Finally $h_p^\Omega$ has smooth strictly positive density relative
to Haar probability, with uniform positive upper and lower bounds
on the compact bundle. The same holds between $L_\Omega\mu$ and
$L\mu$ for any base measure. They have the same supports and the
same pluripolar noncharging property, and neither has point atoms.
Theorem~\ref{thm:classification} finishes the assertions.
\end{proof}

In particular any desired charging fraction is obtained from a
base point mass plus smooth base volume. The conditional density
is the normalized $\Omega$-area; it is Haar for the standard metric.

\section{Scope}
The necessary intermediate-current and closed-semipositive support
and $dd^c$-closed complementary-test restrictions hold for arbitrary
honest shrinking sequences on compact Hermitian manifolds. The complete
support and prescribed-measure classification holds for all such
sequences with every Hermitian metric on the named scalar Hopf
threefold. The construction uses the standard Hopf differential
identities, arbitrarily large
smooth masses and the stipulated shrinking properties as supplied
background. The current structure, averaging, mixed-term controls,
curvature repairs, exact mass corrections and complete weak limits
are proved here. No defect of a particular radial profile or formal
averaging of nonlinear measures is assumed.

The Hopf examples settle both possibilities in the universal
pluripolar-charging question: defects may charge none, or may place
their full excess mass on a pluripolar set. They do not give a
sufficiency theorem for the necessary current conditions on every
Hermitian manifold. The general S01 programme remains open.
\begin{thebibliography}{9}
\bibitem{Demailly}
J.-P. Demailly, \emph{Complex Analytic and Differential Geometry},
version of 21 June 2012, I.4.15--4.24 and I.5.1--5.16.
\end{thebibliography}
\end{document}
